% ============================================================
%  SY02-46 · §2 主定理 · 印刷页 6–8（底本 p6 下段起，p7、p8）
%  条目：Definition 2.1、Theorem 2.2（FGV 主定理）
% ============================================================
\section{主定理}\srcpage{6}

设 $n\ge 2$ 是整数。$\Pic(X)$ 的自然推广，更确切地说是概形 $\Pic(X)$ 在 $\Gm$ 的平凡作用下的商叠的自然推广，是 Riemann 面 $X$ 上秩 $n$ 的向量丛的叠 $\Fib_n(X)$。这个叠是代数的，光滑且纯维数为 $n^2(g-1)$，并且是有限型开集的递增并。若 $g\ge 2$，它含一个稠密开集，该开集是 $X$ 上稳定丛的拟射影模概形上的 $\Gm$-gerbe。$\Fib_n(X)$ 的连通分支 $\Fib_n^d(X)$ 由泛向量丛的次数 $d$ 划分。

\srcpage{7}

$X$ 上秩 $n$ 的向量丛 $\cL$ 也可以看作秩 $n$ 的局部自由 $\cO_X$-模。$\cL$ 在点 $x\in X$ 处的初等下修改，指 $X$ 上秩 $n$ 的向量丛 $\cL'$，它作为子 $\cO_X$-模含于 $\cL$ 中，使得凝聚 $\cO_X$-模 $\cL/\cL'$ 是集中于 $x$、长度为 $1$ 的摩天层。因此 $\cL'$ 的次数等于
\[
  \deg(\cL') = \deg(\cL) - 1.
\]

Hecke 叠是三元组 $(x,\cL,\cL'\subset\cL)$ 的模叠 $\Hck_n(X)$，其中 $x\in X$，$\cL\in\Fib_n(X)$，而 $\cL'\subset\cL$ 是 $\cL$ 在 $x$ 处的初等下修改。

Hecke 对应是图表
\[
  \begin{tikzcd}[column sep=1.8em, row sep=1.8em]
    & \Hck_n(X) \arrow[dl, "{(q,p')}"'] \arrow[dr, "p"] & \\
    X\times\Fib_n(X) & & \Fib_n(X)
  \end{tikzcd}
\]

其中投影 $(q,p')$ 与 $p$ 由
\[
  (q,p')(x,\cL,\cL'\subset\cL) = (x,\cL')
\]
与
\[
  p(x,\cL,\cL'\subset\cL) = \cL
\]
给出，并且二者分别是相对维数 $n-1$ 与 $n$ 的可表、射影、光滑态射。

这个对应在导出范畴 $\Dbc(\Fib_n(X))$ 上诱导一个算子，称为 Hecke 算子。它把 $K\in\Dbc(\Fib_n(X))$ 送到
\[
  H(K) = (q,p')_* p^*K[n-1] \in \Dbc(X\times\Fib_n(X)).
\]

可以「迭代」$H$，特别地得到算子
\[
  H\circ H : \Dbc(\Fib_n(X)) \to \Dbc(X\times X\times\Fib_n(X)),
\]
它在对角线的余集 $(X\times X-\Delta_X)\times\Fib_n(X)$ 上的限制，关于置换 $X$ 的两个副本的作用是 $S_2$-等变的。

\begin{definition}
设 $E$ 是 $X$ 上秩 $n$ 的不可约局部系统。$\Fib_n(X)$ 上的不可约反常层 $K$ 称为相对 $E$ 具有 Hecke 性质，若存在同构
\[
  H(K)\cong E\boxtimes K,
\]
使得诱导同构
\[
  (H\circ H)(K)\cong E\boxtimes E\boxtimes K
\]
在开集 $(X\times X-\Delta_X)\times\Fib_n(X)$ 上的限制是 $S_2$-等变的。
\end{definition}

\srcpage{8}

\begin{theorem}[Drinfeld 于 $n=2$，{[Dr 2]}；一般情形由 Frenkel、Gaitsgory 与 Vilonen 给出，{[F-G-V 2]} 与 {[Ga]}]
对 $X$ 上任一秩 $n$ 的不可约局部系统 $E$，存在 $\Fib_n(X)$ 上的反常层 $\Aut_E$。它到每个连通分支 $\Fib_n^d(X)$ 的限制不可约，并且相对 $E$ 具有 Hecke 性质。
\end{theorem}

定理的证明分两步。

首先，在偶 $(L,s)$ 的模空间 $\Fib'_n(X)$ 上构造可构造层的复形 $\Aut'{}_E$。这里的偶 $(L,s)$ 由 $X$ 上秩 $n$ 的向量丛 $L$ 与一个模数乘的非零截面
\[
  s : (\OmX)^{\otimes(n-1)} \hookrightarrow L
\]
组成。

$\mathrm{GL}(n)$ 尖点自守形式的 Fourier 级数展开，以及关于该级数系数的 Shintani 公式，启发了这个构造；其中还考虑到 Grothendieck 的函数–层字典。但它不是简单的照搬。

然后证明：在 $\Fib'_n(X)$ 的一个大的开集上，$\Aut'{}_E$ 是反常层；它沿遗忘截面 $s$ 的态射 $\Fib'_n(X)\to\Fib_n(X)$ 下降，成为 $\Fib_n(X)$ 上具有所需性质的反常层 $\Aut_E$。当然，全部困难都在第二步。
